\section{The Language of Predicate Logic}
\begin{questions}
    % An Exercise in Languages
    \item \begin{partquestions}{\alph*}
        % \item Some possible strings are ``12121'', ``12312'', ``12321'', ``23123'', and ``32323''.
        \item The alphabet of $L$ must contain the symbols $a$, $b$, $c$, $d$, and $r$.
        \item {\L}ukasiewicz notation: ``$+\;4\;\times\;-\;5\;6\;7$''. Reversed {\L}ukasiewicz notation: ``$4\;5\;6\;-\;7\;\times\;+$''.
        \item The formal grammar of $L$ contains only one formation rule.
        \begin{quote}
            If $s \in \Alphabet$ is $n$-ary and $\phi_1, \phi_2, \phi_3, \dots, \phi_n$ are words in $L$, then
            \[
                \phi_1\phi_2\phi_3\dots\phi_ns
            \]
            is also a word in $L$.
        \end{quote}
        This is similar to {\L}ukasiewicz notation, except that the $s$ is placed at the end instead of at the front.
    \end{partquestions}

    % Language Construction Puzzle
    \item To argue which of the strings are words and which are not, observe that we may only have at most one $a$ that may come after a $b$, regardless of the number of $b$s that are in between.
    \begin{partquestions}{\alph*}
        \item Word. $ba \overset{1}{\longrightarrow} aba \overset{2}{\longrightarrow} abab$.
        \item Word. $\emptystring \overset{1}{\longrightarrow} a \overset{1}{\longrightarrow} aa \overset{1}{\longrightarrow} aaa \overset{2}{\longrightarrow} aaab \overset{2}{\longrightarrow} aaabb \overset{2}{\longrightarrow} aaabbb$.
        \item Not a word by above observation. From the 1st $b$ there are 2 $a$s that come after it.
        \item Word. $ba \overset{1}{\longrightarrow} aba \overset{2}{\longrightarrow} abab \overset{2}{\longrightarrow} ababb \overset{2}{\longrightarrow} ababbb$.
    \end{partquestions}

    % Strings in $\SetLang$
    \item All of them are strings, but most are not words.
    
    % What's in $\SetLang$?
    \item \begin{partquestions}{\alph*}
        \item This is a term. From \textbf{T0} we know that $x$ is a term, and \textbf{T1} tells us that 7, being a nullary function symbol, is also a term. Thus \textbf{T1} tells us that $f(x,7)$ is a term.
        \item This is a \wff, but not an atomic formula since quantifiers are being used in the string.
        \item Neither a term nor a \wff. Observe that both $f(x)$ and 3 are terms by \textbf{T1}, but there is no rule (for terms and for \wffs) that allows concatenating the two terms together into the form $3(f(x))$.
        \item This is a \wff, but not an atomic formula.
        \item This is a \wff, but not an atomic formula.
        \item Neither a term nor a \wff. The string $x=y$ is a \wff~by \textbf{F1}, so it is not a term. Hence having ``$x=y$'' being used as an argument for the predicate symbol $P$ is invalid.
    \end{partquestions}

    % Bound and Free Variables
    \item \begin{partquestions}{\alph*}
        \item Bound.
        \item Both bound and free. The left side (i.e., ``$x=f(y)$'') has $x$ being a free variable, while the right side (i.e., ``$\forall x \exists y (y = x)$'') has $x$ being a bound variable.
        \item Free.
        \item Neither. The variable $x$ does not appear in the \wff.
    \end{partquestions}
\end{questions}
